\documentclass[11pt, a4paper]{article}

% ── Packages ──
\usepackage[utf8]{inputenc}
\usepackage[T1]{fontenc}
\usepackage{lmodern}
\usepackage[margin=2.5cm]{geometry}
\usepackage{amsmath, amssymb, amsfonts}
\usepackage{graphicx}
\usepackage{booktabs}
\usepackage{hyperref}
\usepackage{xcolor}
\usepackage{listings}
\usepackage{float}
\usepackage{enumitem}
\usepackage{fancyhdr}
\usepackage{tabularx}
\usepackage{caption}
\usepackage{algorithm}
\usepackage{algpseudocode}

% ── Colors ──
\definecolor{codeblue}{rgb}{0.1, 0.2, 0.6}
\definecolor{codegray}{rgb}{0.5, 0.5, 0.5}
\definecolor{codebg}{rgb}{0.96, 0.97, 0.98}

% ── Listings ──
\lstset{
  backgroundcolor=\color{codebg},
  basicstyle=\ttfamily\small,
  keywordstyle=\color{codeblue}\bfseries,
  commentstyle=\color{codegray}\itshape,
  numbers=left, numberstyle=\tiny\color{codegray},
  breaklines=true, frame=single, rulecolor=\color{codegray!40},
  tabsize=4, showstringspaces=false
}

% ── Header/Footer ──
\pagestyle{fancy}
\fancyhf{}
\fancyhead[L]{\small National Space Hackathon 2026}
\fancyhead[R]{\small Technical Report — Project AETHER}
\fancyfoot[C]{\thepage}
\renewcommand{\headrulewidth}{0.4pt}

\hypersetup{colorlinks=true, linkcolor=codeblue, urlcolor=codeblue, citecolor=codeblue}

% ── Title ──
\title{
  \vspace{-1cm}
  \textbf{Project AETHER} \\[4pt]
  \large Autonomous Constellation Manager for \\
  Orbital Debris Avoidance and Station-Keeping \\[8pt]
  \normalsize National Space Hackathon 2026 — Technical Report
}
\author{
  Pranjal Saxena
}
\date{March 2026}

\begin{document}
\maketitle
\tableofcontents
\newpage

% ══════════════════════════════════════════════════════════
\section{Executive Summary}
% ══════════════════════════════════════════════════════════

Project AETHER (\textbf{A}utonomous \textbf{E}vasion, \textbf{T}racking, \textbf{H}euristic \textbf{E}ngine for \textbf{R}esilient orbits) is a full-stack Autonomous Constellation Manager (ACM) that manages a fleet of 50+ LEO satellites navigating a hazardous debris environment of 10{,}000+ tracked objects. The system ingests high-frequency orbital telemetry, predicts conjunctions up to 24 hours ahead using KD-tree spatial indexing, autonomously computes and executes collision-avoidance and station-keeping maneuvers, and presents real-time situational awareness through a 3D operational dashboard.

Key technical achievements:
\begin{itemize}[nosep]
  \item J2-perturbed orbital propagation via 4th-order Runge-Kutta with Numba JIT — 10{,}000 objects propagated in $\sim$0.1\,s
  \item $O(N \log N)$ conjunction screening via \texttt{scipy.spatial.KDTree} with coarse-then-fine TCA refinement
  \item Constraint-validated maneuver scheduling enforcing cooldown, LOS, signal delay, fuel budget, and max-$\Delta v$ limits
  \item Blind-conjunction handling: pre-uploading evasion sequences before satellite enters communication blackout
  \item Three.js 3D globe visualization rendering 10{,}000+ debris at 60\,FPS via WebGL \texttt{InstancedMesh}
\end{itemize}

% ══════════════════════════════════════════════════════════
\section{System Architecture}
% ══════════════════════════════════════════════════════════

\subsection{Overview}
The ACM is a containerized application (Docker, \texttt{ubuntu:22.04}) exposing a RESTful API on port~8000. It comprises a Python/FastAPI backend for physics computation and autonomous decision-making, and a React/Three.js frontend for visualization.

\subsection{Module Decomposition}

\begin{table}[H]
\centering
\caption{Backend module responsibilities}
\begin{tabularx}{\textwidth}{l X}
\toprule
\textbf{Module} & \textbf{Responsibility} \\
\midrule
\texttt{physics/propagator.py} & RK4 numerical integrator with J2 perturbation; batch-parallel via Numba \\
\texttt{physics/coordinates.py} & ECI$\leftrightarrow$ECEF$\leftrightarrow$Geodetic transformations; GMST computation \\
\texttt{physics/maneuver.py} & RTN frame construction; evasion/recovery $\Delta v$ planning; Tsiolkovsky fuel equation \\
\texttt{engine/state\_manager.py} & In-memory state store using NumPy arrays for vectorized operations \\
\texttt{engine/conjunction.py} & KD-tree spatial screening + coarse/fine TCA search \\
\texttt{engine/scheduler.py} & Constraint-validated maneuver queue; auto-evasion; EOL management \\
\texttt{engine/ground\_stations.py} & Line-of-sight computation; visibility window prediction \\
\texttt{engine/station\_keeping.py} & Nominal slot drift monitoring; uptime tracking; recovery scheduling \\
\texttt{utils/logger.py} & Structured JSON event logging for all maneuvers and detections \\
\bottomrule
\end{tabularx}
\end{table}

\subsection{Data Flow}

\begin{enumerate}[nosep]
  \item Grader sends object states via \texttt{POST /api/telemetry}
  \item State manager ingests and indexes all objects
  \item Conjunction detector screens pairs via KD-tree, refines close approaches
  \item Scheduler auto-queues evasion burns for critical CDMs, validates constraints
  \item Grader advances time via \texttt{POST /api/simulate/step}
  \item Propagator integrates all orbits; scheduler executes due burns
  \item Station-keeping manager checks slot compliance, schedules recovery burns
  \item Frontend polls \texttt{GET /api/visualization/snapshot} for real-time display
\end{enumerate}

% ══════════════════════════════════════════════════════════
\section{Numerical Methods}
% ══════════════════════════════════════════════════════════

\subsection{Equations of Motion (J2 Perturbation)}

All propagation uses the J2-perturbed two-body equations of motion in the ECI~(J2000) frame as specified in the problem statement:

\begin{equation}
  \ddot{\vec{r}} = -\frac{\mu}{|\vec{r}|^3}\,\vec{r} + \vec{a}_{J_2}
\end{equation}

where $\mu = 398600.4418\;\text{km}^3/\text{s}^2$ and the J2 acceleration is:

\begin{equation}
  \vec{a}_{J_2} = \frac{3}{2}\frac{J_2\,\mu\,R_E^2}{|\vec{r}|^5}
  \begin{pmatrix}
    x\left(\frac{5z^2}{|\vec{r}|^2} - 1\right) \\[4pt]
    y\left(\frac{5z^2}{|\vec{r}|^2} - 1\right) \\[4pt]
    z\left(\frac{5z^2}{|\vec{r}|^2} - 3\right)
  \end{pmatrix}
\end{equation}

with $R_E = 6378.137\;\text{km}$ and $J_2 = 1.08263 \times 10^{-3}$.

\subsection{4th-Order Runge-Kutta Integration}

The state vector $\vec{y} = [\vec{r},\, \vec{v}]^T$ is integrated using the classical RK4 scheme:

\begin{equation}
  \vec{y}_{n+1} = \vec{y}_n + \frac{h}{6}\left(\vec{k}_1 + 2\vec{k}_2 + 2\vec{k}_3 + \vec{k}_4\right)
\end{equation}

with the standard evaluation points. Our default timestep is $h = 10$\,s, chosen to balance accuracy against computational cost for LEO orbits with periods of $\sim$90\,min.

\subsection{Performance Optimization}

All propagation kernels use \textbf{Numba JIT} compilation (\texttt{@njit}) to achieve near-C performance. Batch propagation uses \texttt{@njit(parallel=True)} with \texttt{prange} to distribute across CPU cores. This achieves:

\begin{itemize}[nosep]
  \item Single object, 1 orbit ($\sim$5400\,s): $\sim$0.5\,ms
  \item 10{,}000 objects, 3600\,s step: $\sim$100\,ms on 8 cores
\end{itemize}

% ══════════════════════════════════════════════════════════
\section{Spatial Optimization: Conjunction Detection}
% ══════════════════════════════════════════════════════════

\subsection{The $O(N^2)$ Problem}

Na\"{\i}vely checking every satellite--debris pair costs $O(N_s \cdot N_d)$ per timestep. With $N_s = 50$ and $N_d = 10{,}000$, this is 500{,}000 distance computations per step, multiplied across hundreds of propagation steps for a 24-hour prediction horizon.

\subsection{KD-Tree Screening ($O(N \log N)$)}

We construct a \texttt{scipy.spatial.KDTree} over the debris positions and query a screening ball of radius $R_\text{screen} = 500$\,km around each satellite. This reduces the candidate set from $O(N_d)$ to $O(k)$ where $k \ll N_d$ for each satellite.

\subsection{Coarse-Then-Fine TCA Search}

For each candidate pair surviving the KD-tree screen:
\begin{enumerate}[nosep]
  \item \textbf{Coarse pass}: Propagate both objects at 60\,s intervals over 24 hours; record the time of minimum distance.
  \item \textbf{Fine pass}: Re-propagate a $\pm$120\,s window around the coarse minimum at 1\,s intervals to refine the Time of Closest Approach (TCA) and miss distance.
\end{enumerate}

This two-phase approach avoids expensive fine-grained propagation for the vast majority of pairs that are not true threats.

\begin{algorithm}[H]
\caption{Conjunction Assessment Pipeline}
\begin{algorithmic}[1]
\State $\text{tree} \gets \text{KDTree}(\text{debris\_positions})$
\For{each satellite $s$}
  \State $\text{nearby} \gets \text{tree.query\_ball\_point}(s.\vec{r},\; R_\text{screen})$
  \For{each debris $d$ in nearby}
    \State $(t_\text{min}, d_\text{min}) \gets \text{coarse\_search}(s, d, \Delta t = 60\text{s}, T = 86400\text{s})$
    \If{$d_\text{min} < 5\;\text{km}$}
      \State $(t_\text{TCA}, d_\text{TCA}) \gets \text{fine\_search}(s, d, t_\text{min} \pm 120\text{s}, \Delta t = 1\text{s})$
      \State Classify risk: CRITICAL ($<$100\,m), RED ($<$1\,km), YELLOW ($<$5\,km)
    \EndIf
  \EndFor
\EndFor
\end{algorithmic}
\end{algorithm}

% ══════════════════════════════════════════════════════════
\section{Maneuver Planning and Execution}
% ══════════════════════════════════════════════════════════

\subsection{RTN Frame and Evasion Strategy}

Maneuvers are computed in the satellite's local Radial-Transverse-Normal (RTN) frame:

\begin{equation}
  \hat{R} = \frac{\vec{r}}{|\vec{r}|}, \quad
  \hat{N} = \frac{\vec{r} \times \vec{v}}{|\vec{r} \times \vec{v}|}, \quad
  \hat{T} = \hat{N} \times \hat{R}
\end{equation}

Our primary evasion strategy applies a \textbf{transverse burn} ($\hat{T}$ direction), which is the most fuel-efficient way to shift the satellite's along-track timing. The required $\Delta v$ scales inversely with lead time, motivating early detection and scheduling.

\subsection{Fuel Mass Depletion (Tsiolkovsky)}

Every burn depletes propellant according to the rocket equation:

\begin{equation}
  \Delta m = m_\text{current}\left(1 - e^{-\frac{|\Delta\vec{v}|}{I_{sp}\,g_0}}\right)
\end{equation}

with $I_{sp} = 300$\,s and $g_0 = 9.80665$\,m/s$^2$. The system dynamically tracks mass, meaning later burns on a lighter satellite are slightly more efficient.

\subsection{Constraint Validation}

Every maneuver command passes through 6 validation gates before acceptance:

\begin{table}[H]
\centering
\begin{tabularx}{\textwidth}{l X}
\toprule
\textbf{Constraint} & \textbf{Implementation} \\
\midrule
Max $\Delta v$ & $|\Delta\vec{v}| \leq 15$\,m/s per burn \\
Signal delay & Burn time $\geq$ current time $+ 10$\,s \\
Fuel budget & Tsiolkovsky mass check against remaining propellant \\
Cooldown & 600\,s minimum between burns on the same satellite \\
Ground station LOS & Elevation angle $\geq$ station mask at upload time \\
Satellite existence & Valid ID check in state registry \\
\bottomrule
\end{tabularx}
\end{table}

\subsection{Blind Conjunction Handling}

If a critical conjunction is predicted while the satellite is in a communication blackout zone:
\begin{enumerate}[nosep]
  \item The system propagates the satellite forward to find the next ground station contact window.
  \item If the contact window occurs \emph{before} the TCA, the evasion sequence is scheduled for upload during that window.
  \item If no contact window exists before TCA, the maneuver cannot be executed (logged as unresolvable).
\end{enumerate}

\subsection{Station-Keeping and Recovery}

After evasion, the system monitors drift from the nominal orbital slot (10\,km tolerance). When drift is detected and no conflicting maneuver is pending, a phasing recovery burn is automatically scheduled in the prograde/retrograde direction to adjust the orbital period and allow the satellite to drift back.

\subsection{End-of-Life Management}

When fuel drops below 5\% ($\leq 2.5$\,kg), a graveyard orbit raise of $\sim$25\,km is autonomously scheduled to prevent the satellite from becoming uncontrollable debris.

% ══════════════════════════════════════════════════════════
\section{Ground Station Network and Communication}
% ══════════════════════════════════════════════════════════

\subsection{Line-of-Sight Geometry}

For each of the 6 ground stations, the elevation angle of a satellite is computed as:

\begin{equation}
  \sin(\varepsilon) = \frac{(\vec{r}_\text{sat} - \vec{r}_\text{gs}) \cdot \hat{u}_\text{gs}}{|\vec{r}_\text{sat} - \vec{r}_\text{gs}|}
\end{equation}

where $\hat{u}_\text{gs}$ is the local ``up'' unit vector at the ground station (its ECEF position normalized). Visibility requires $\varepsilon \geq \varepsilon_\text{min}$ (station-specific mask angle from 5$^\circ$ to 15$^\circ$).

This formulation naturally accounts for Earth curvature: a satellite behind Earth yields a negative or sub-threshold elevation.

\subsection{ECI to ECEF Conversion}

The ECI-to-ECEF rotation uses Greenwich Mean Sidereal Time:
\begin{equation}
  \theta_\text{GMST} = 280.46061837^\circ + 360.98564736629 \cdot (JD - 2451545.0) + \ldots
\end{equation}

applied as a $z$-axis rotation to transform satellite positions from the inertial to Earth-fixed frame for ground station geometry.

% ══════════════════════════════════════════════════════════
\section{Frontend: Orbital Insight Dashboard}
% ══════════════════════════════════════════════════════════

The frontend is a single-page React application with four PS-required visualization modules:

\begin{enumerate}[nosep]
  \item \textbf{3D Globe (Ground Track View)}: Three.js Earth sphere with NASA Blue Marble texture, satellite markers, debris cloud via \texttt{InstancedMesh} ($\sim$1 draw call for 10K points), ground station markers, and threat lines. Supports orbit controls (rotate, zoom).

  \item \textbf{Conjunction Bullseye Plot}: SVG polar chart with the selected satellite at center. Radial distance encodes TCA; angle encodes approach vector. Color indicates risk level (red/amber/green).

  \item \textbf{Fleet Fuel Heatmap}: 10$\times$5 grid of color-coded cells per satellite showing propellant percentage, plus a bar chart and min/avg/max summary.

  \item \textbf{Maneuver Timeline (Gantt)}: Horizontal tracks per satellite showing evasion burns (red), recovery burns (cyan), and 600\,s cooldown blocks (hatched), with a ``NOW'' marker.
\end{enumerate}

\textbf{Performance}: The \texttt{InstancedMesh} approach renders 10{,}000 debris objects as a single GPU draw call. Combined with React Three Fiber's reconciler, the dashboard maintains 60\,FPS with full interactivity.

% ══════════════════════════════════════════════════════════
\section{Technology Stack}
% ══════════════════════════════════════════════════════════

\begin{table}[H]
\centering
\begin{tabularx}{\textwidth}{l l X}
\toprule
\textbf{Layer} & \textbf{Technology} & \textbf{Purpose} \\
\midrule
API Framework & FastAPI + Uvicorn & Async REST API on port 8000 \\
Physics & NumPy + Numba & Vectorized propagation with JIT compilation \\
Spatial Index & SciPy KDTree & $O(N\log N)$ conjunction screening \\
Frontend & React 18 + Vite & Component-based dashboard with HMR \\
3D Rendering & Three.js + R3F & WebGL Earth globe and orbital visualization \\
Charts & SVG + Recharts & Polar plots, fuel gauges, timelines \\
Containerization & Docker (ubuntu:22.04) & Reproducible deployment \\
Logging & JSON structured logs & Machine-readable event stream \\
\bottomrule
\end{tabularx}
\end{table}

% ══════════════════════════════════════════════════════════
\section{Evaluation Criteria Mapping}
% ══════════════════════════════════════════════════════════

\begin{table}[H]
\centering
\begin{tabularx}{\textwidth}{l c X}
\toprule
\textbf{Criterion} & \textbf{Weight} & \textbf{Our Approach} \\
\midrule
Safety Score & 25\% & Autonomous CRITICAL CDM evasion; 2\,km safety margin; blind-conjunction pre-upload \\
Fuel Efficiency & 20\% & Transverse-first strategy; early detection reduces $\Delta v$; mass-aware Tsiolkovsky \\
Constellation Uptime & 15\% & Automatic recovery burns; drift monitoring; exponential penalty tracking \\
Algorithmic Speed & 15\% & KD-tree $O(N\log N)$; Numba JIT parallel propagation; coarse/fine TCA \\
UI/UX \& Visualization & 15\% & 3D WebGL globe; 60\,FPS InstancedMesh; all 4 required modules \\
Code Quality \& Logging & 10\% & Modular architecture; structured JSON logs; type-validated API \\
\bottomrule
\end{tabularx}
\end{table}

% ══════════════════════════════════════════════════════════
\section{Conclusion}
% ══════════════════════════════════════════════════════════

Project AETHER demonstrates that autonomous constellation management is achievable with efficient numerical methods and intelligent constraint handling. Our KD-tree spatial indexing eliminates the $O(N^2)$ bottleneck, Numba-accelerated propagation handles 10{,}000+ objects in real-time, and the six-gate maneuver validation pipeline ensures no burn violates the physical or operational constraints of the mission. The 3D operational dashboard provides the situational awareness needed for human-in-the-loop oversight while the autonomous engine handles the high-frequency decision-making that manual operations cannot scale to.

\end{document}
